\section{Goal, motivation, and results}
\label{sec:context}

\subsection{Source-defined language-model inference}

LeanExe's goal is to run an algorithm written in Lean with a checked proof that the emitted program implements that algorithm.  The source definition specifies the intended computation.  The compiled program supplies an execution path independent of Lean's runtime.  A per-artifact proof connects the target instructions, their memory operations, and their observable results to the chosen specification.  GPT-2 provides a case in which that connection must survive millions of floating-point operations, nested tensor loops, retained key/value state, and repeated allocation and release~\cite{subset,cached,entry}.

The initial practical objective was command-line inference with a pretrained language model and a complete source-agreement proof.  A user supplies a prompt, the host converts it into vocabulary indices, and the inference module returns the scores used to select subsequent tokens.  The implemented model has the block structure and parameter dimensions of GPT-2 124M, with an implemented context limit of 128 tokens.  Its weights arrive as runtime bytes.  The verification therefore applies to replacement weight values of the same shape without retraining or checking a new weight-specific theorem~\cite{manifest,guide,session}.

The algorithm fixes FP32 operation order and the numerical routines for normalization, exponentiation, and activation.  Equality with this source algorithm requires equality of represented words, including the specified exceptional results.  The verification target also includes successful execution: a proof of output equality conditional on a completed call would leave allocation failure and nontermination unresolved.  The session theorem combines these obligations and derives the memory conditions used by each token call from initialization~\cite{entry,session,budget}.

GPU execution introduces a second implementation of selected operations.  Matrix-product output coordinates can execute independently, while each coordinate retains its ordered accumulation.  The WGSL development moves those products to shaders and keeps the rest of the model in WebAssembly.  The resulting proof must connect three interfaces: the Lean kernel definition to emitted shader statements, the shader's input and output views to the packed model tensors, and the external call to the WebAssembly controller.  The conditional hybrid theorem composes these connections~\cite{wg-plan,wg-packedsource,wg-backend}.

\subsection{Two execution paths and their checked results}

The CPU artifact result starts with the complete frozen WebAssembly byte sequence.  Kernel-checked decoding and validation recover a module whose Talos translation equals the module used by the inference proof.  Talos is a WebAssembly interpreter and proof library written in Lean.  In its semantics, the recovered module terminates and returns the exact cache and logits of the Lean recurrence for every permitted session.  The public theorem is \longcode{Project.Gpt2CachedStep.Artifact.artifact\_gpt2\_128\_exact}~\cite{artifact}.

Wasmtime is the standalone WebAssembly runtime used for native CPU execution.  Its selected configuration canonicalizes arithmetic NaNs to agree with the formal word-level computation.  Application of the theorem to the running command trusts Wasmtime, the native host, and the association between file bytes and the embedded binary.  The recorded artifact check compares those bytes.  Section~\ref{sec:boundary} states the remaining runtime and automation assumptions~\cite{wasmtimeintro,wasmtime,artifactgate}.

The hybrid result starts with a modeled WebAssembly controller and six certified shader strings.  WGSL is the WebGPU Shading Language.  WebGPU supplies the API through which a host creates shader pipelines, submits work, and transfers buffers.  The checked shader results concern an explicit statement semantics for the supported grammar.  The complete hybrid theorem assumes that external calls complete according to that strict shader semantics, copy the specified bytes, and preserve the rest of the WebAssembly store.  Packed-layout theorems then identify those results with the CPU source operations.  The public hybrid session declarations include \code{gpt2\_128\_exact} and \code{gpt2\_128\_exact\_for} in namespace \code{Project.Gpt2Hybrid.Spec}~\cite{wg-execution,wg-backend,wg-composition}.

These two statements support different deployment claims.  The CPU proof includes the binary-to-model connection.  The hybrid proof includes a checked shader-to-source connection and an explicit external-execution premise.  The current hybrid binary enters the proof through declarations generated by an external parser.  Its native and browser hosts provide implementations of the required call sequence, and execution tests provide evidence about those implementations.  The theorem's explicit host premise and that empirical evidence remain separate~\cite{wg-build,wg-host,wg-review}.

\subsection{Execution evidence and report checkpoints}

The CPU test record compares 6,432,896 logits against PyTorch over one 128-token sequence and all its prefixes.  Every component meets the recorded tolerance, and the maximum observed absolute difference is approximately 0.001450.  The hybrid integration journal records the same number of exact logit comparisons against the parent WebAssembly path on a CPU WebGPU implementation.  The browser screenshots supplied for this report show both backends completing the same 32-token continuation.  The page reports a decode-throughput ratio of 15.72.  Sections~\ref{sec:tests} and~\ref{sec:browser} identify the populations, measurement rules, and provenance of these observations~\cite{tests,wg-journal}.

The CPU proof checkpoint is \longcode{4360920c3060d1c625b1860229b5116804d177c8}.  The WGSL and browser source checkpoint is \longcode{9c7c7898ecae5f636cf1047142c8a68c1936e041}.  Its two commits after the earlier WGSL report add the comparison page and focused worker tests.  They change host, interface, and documentation files.  The shader and inference proof sources retain the preceding reviewed development.  Both paths use Lean 4.34.0-rc2 at revision \longcode{6a10ac8c22beadecabdbb0919c2b50214762f91d} and Talos at revision \longcode{87e3aa5e8f6e6f3b3eb5e7e4c5aba43071002d47}~\cite{repo,wg-checkpoint,talos}.

This report combines and expands the CPU and WGSL technical reports~\cite{cpu-report,wg-report}.  It states the algorithm, proof subjects, execution assumptions, and measurements in one account.  The literature search compares complete artifact verification with source-level network reasoning, compiler transformations, and cryptographic certificates for inference evaluations.  The search found no earlier result with the same combination of a complete GPT-style inference module, proof-assistant-checked source agreement, and a checked connection from executable binary bytes.  The priority assessment is qualified by the search coverage described in Section~\ref{sec:related}.
